% =====================================================================
% Plastic and Reconstructive Surgery (PRS): Original Article WORKING SHELL
% Unofficial manuscript class, NOT the journal's production template.
% Build: pdflatex main; biber main; pdflatex main; pdflatex main
% Toggle [draftfigures] below for placeholders; remove for submission.
% Original Articles: structured abstract <=250 words, body <=3000 words.
% Figure pieces + tables combined: <=20 (verify current article category).
% =====================================================================
% Submission: use [submission] or [submission,draftfigures]
% Visual approximation of a published paper: use [journalpreview,draftfigures]
% \documentclass[submission,draftfigures]{prsmanuscript}
\documentclass[journalpreview,draftfigures]{prsmanuscript}
\usepackage[backend=biber,style=numeric-comp,sorting=none,
  maxbibnames=6,minbibnames=3,giveninits=true,
  doi=true,isbn=false,url=false]{biblatex}
\addbibresource{references.bib}
% Superscript numbers after punctuation, e.g. a sentence.\supercite{key}
% AMA-like numeric reference list; verify punctuation/journal abbreviations
% against the journal's current examples before final submission.
\DeclareFieldFormat{labelnumberwidth}{#1\adddot}
\DeclareFieldFormat[article]{title}{#1}
\DeclareFieldFormat{journaltitle}{\mkbibemph{#1}}
\renewbibmacro{in:}{}

% ----- Manuscript metadata -----
\PRStitle{Placeholder Title: A Prospective Evaluation of a Reconstructive Technique}
\PRSshorttitle{Evaluation of a Reconstructive Technique}
\PRSauthors{First A. Author, MD\textsuperscript{1}; Second B. Author, MSc\textsuperscript{2}; Third C. Author, MD\textsuperscript{1}}
\PRSaffiliations{\textsuperscript{1}Department of Plastic Surgery, Example Institution, City, State/Province, Country\\
\textsuperscript{2}Department of Biomedical Engineering, Example University, City, State/Province, Country}
\PRScorrespondence{First A. Author, Department, Institution, mailing address; email@example.org}
\PRSwordcounts{[insert actual count]}{[insert actual count]}
\PRSmetadata{\noindent\textbf{Article category:} Original Article\par
\noindent\textbf{Prior presentation:} [Meeting, location, date; or None]\par
\noindent\textbf{Funding:} [Funding source and grant; or None]\par
\noindent\textbf{Conflicts of interest:} [Disclosures; or None]\par}

\begin{document}
\PRSmaketitle

% For an Experimental Article, also add a 1--2 sentence Clinical Relevance
% Statement at the END of the abstract and in the submission field.
\begin{PRSabstract}
\abstractpart{Background}{[Two or three sentences describing the problem, evidence gap, and objective.]}
\abstractpart{Methods}{[Study design, setting, participants/samples, intervention or technique, primary outcome, and analysis.]}
\abstractpart{Results}{[Sample size; effect size; uncertainty interval; key comparison and numerical findings. Do not invent results.]}
\abstractpart{Conclusions}{[A concise, evidence-supported answer to the study objective.]}
\PRSkeywords{reconstructive surgical procedures; surgical techniques; outcome assessment; [add appropriate MeSH terms]}
\end{PRSabstract}
\PRSstartbody % Enables two-column body ONLY in journalpreview mode.

\section{Introduction}
[Introduce the clinical problem and the current standard of care. Then describe the relevant knowledge gap.]
A representative numbered citation appears at the end of this sentence.\supercite{exampletrial}
A second publication can be cited with a combined citation.\supercite{exampletrial,examplereview}

[End with a specific study objective and, where relevant, a prespecified hypothesis.]

\section{Methods}
\subsection{Study Design and Ethics}
[Design, institution/setting, dates, institutional review board approval or exemption, informed consent, registration (when relevant).]

\subsection{Participants and Procedures}
[Inclusion and exclusion criteria, recruitment, intervention or device details, comparator, and reproducibility information.]
Use \emph{italic text} for emphasized terminology, \textbf{bold} sparingly, and an explanatory footnote when necessary.\footnote{Footnotes should be used sparingly; essential methods belong in the main text.}

\subsection{Outcome Measures and Statistical Analysis}
[Define primary and secondary outcomes, sample size rationale, statistical tests, multiplicity handling, missing data, and software.]
An inline equation can be written as $\Delta = y_{\mathrm{post}}-y_{\mathrm{pre}}$.
An optional displayed equation is
\begin{equation}
 \bar{x}=\frac{1}{n}\sum_{i=1}^{n}x_i .\label{eq:mean}
\end{equation}
Equation~\eqref{eq:mean} illustrates automatic equation referencing.

\section{Results}
\subsection{Study Population}
[Include screened/enrolled/analyzed counts and baseline characteristics.]
Refer to Table~\ref{tab:baseline} in numerical order.

\subsection{Primary and Secondary Outcomes}
[Present estimates with measures of precision, denominators, and clinically meaningful interpretation.]
Figure~\ref{fig:overview} is a draft placeholder; the actual figure file must be uploaded separately.
\begin{figure}[tbp]\centering
\PRSfigureplaceholder{1}{Example workflow or outcome chart}
\caption{Placeholder caption for drafting; final legend is collected after the references.}\label{fig:overview}
\end{figure}
% Optional *drafting-only* embedded figure (replace the placeholder above):
% \begin{figure}[htbp]\centering
% \includegraphics[width=.8\linewidth]{figures/figure1.pdf}
% \caption{Draft caption for internal review only.}\label{fig:alt}
% \end{figure}

\section{Discussion}
[State principal findings, situate results among existing studies, explain clinical significance, and discuss plausible explanations.]

\subsection{Limitations}
[Discuss bias, statistical uncertainty, external validity, measurement issues, and generalizability.]

\subsection{Conclusions}
[Close with a short conclusion without introducing new results.]

% ----- Statements: replace rather than retaining untrue stock declarations -----
\section*{Acknowledgments and Declarations}
\PRSdisclosure{Acknowledgments}{[Non-author contributions, if any.]}
\PRSdisclosure{Funding}{[Grant and sponsor roles; or no external funding.]}
\PRSdisclosure{Disclosure}{[Relevant financial/nonfinancial interests for all authors; or none.]}
\PRSdisclosure{Ethics approval and consent}{[IRB name/approval number, consent arrangements, or legitimate exemption.]}
\PRSdisclosure{Data availability}{[Availability and privacy/access restrictions, if applicable.]}
\PRSdisclosure{Author contributions}{[Contributions; verify all named authors meet authorship criteria.]}

\clearpage
\printbibliography[title={References},heading=bibintoc]

% PRS figure legends: text only; images separately uploaded.
\clearpage
\section*{Figure Legends}
\PRSfigurelegend{1}{[Description of schematic, imaging modality and view, labels/panels (A, B, etc.), abbreviation definitions, and permissions if adapted.]}

% Tables on separate page(s), not embedded in body; no vertical rules.
\clearpage
\section*{Tables}
\PRStable{tab:baseline}{Example baseline characteristics (replace values with actual data)}
\begin{threeparttable}
\begin{tabularx}{\linewidth}{@{}Xcc@{}}
\toprule
Characteristic & Group A ($n=\cdots$) & Group B ($n=\cdots$) \\
\midrule
Sample size, $n$ & [n] & [n] \\
Age, mean (SD), years & [value] & [value] \\
Outcome, median (IQR) & [value] & [value] \\
\bottomrule
\end{tabularx}
\begin{tablenotes}[flushleft]\footnotesize
\item Abbreviations: IQR, interquartile range; SD, standard deviation.
\item Values shown are placeholders; use appropriate summary statistics for the data.
\end{tablenotes}
\end{threeparttable}

% SDC may be cited in text as Supplemental Digital Content 1; upload separately.
% Cover letter, disclosures, consent/permissions and high-resolution figure
% files generally need separate editorial-system submission.
\end{document}
